\section{Phase-Purity Bottlenecks and Characterization}

The prior experimental conclusions are best read as condition--result pairs rather than as a ranking of methods. Table~\ref{tab:prior} collects representative observations from the corpus. Temperatures or phase names are included only when stated in the cited record; a neighbor system is labeled as such.

\begin{table}[t]
\centering
\caption{Representative prior experimental conclusions. “Direct” denotes the named composition; “neighbor” denotes a bounded analogue.}
\label{tab:prior}
\begin{tabular}{P{0.23\linewidth}P{0.25\linewidth}P{0.37\linewidth}P{0.10\linewidth}}
\toprule
System or route & Reported condition & Result relevant to phase purity or transfer & Relation \\
\midrule
BaY$_2$Si$_3$O$_{10}$ & Molybdate flux growth & YO$_6$ chains and trisilicate groups were resolved in a single-crystal structure & Neighbor \cite{kolitsch2006bay2si3o10} \\
Ba$_2$ZnSi$_2$O$_7$ & Crystal-structure determination & Isolated silicate units define a Ba--Zn structural motif & Neighbor \cite{kaiser2002crystal} \\
Ba$_3$Zn$_4$Si$_4$O$_{15}$ & Structure refinement & Isolated SiO$_4$ and Si$_2$O$_7$ units were identified & Neighbor \cite{ababaikeri2024ba} \\
Ba--Y--Zn--Si--O glass & Multicomponent melt and spectroscopy & Four cations were present, but a phase-pure quaternary crystal was not established & Direct composition \cite{kozhevnikova2024synthesis} \\
BaO--CaO--SiO$_2$ glass-ceramic & Controlled crystallization & Phase fractions changed with heat-treatment conditions & Transfer \cite{hu2025crystallization} \\
ZnO--SiO$_2$ ternaries & Equilibration, quench, EPMA and XRD & Solid solutions and liquidus fields were measured for Zn-bearing silicates & Thermodynamic neighbor \cite{bonner2024phase,jeon2025phase} \\
Ba--Sr--Zn--Si--O glass & Surface and volume crystallization & Growth velocity depended on the crystallization geometry and thermal history & Kinetic neighbor \cite{waurischk2020crystal} \\
\bottomrule
\end{tabular}
\end{table}

\subsection{Diffraction and compositional mapping}

Powder XRD is indispensable for phase indexing, yet a single pattern can miss a minor phase, amorphous fraction, or overlapping reflections. Phase-equilibrium studies therefore combine rapid quenching with electron-probe microanalysis and XRD, allowing phase compositions and assemblages to be evaluated together \cite{kim2025key,derensy2026phase,abdeyazdan2024phase,abdeyazdan2023experimental}. In a BYZSO campaign, bulk XRD should be paired with EPMA or wavelength-dispersive mapping across multiple grains so that Ba/Y/Zn segregation is not averaged away.

\subsection{In-situ and complementary methods}

Hot-stage microscopy and time-resolved scattering reveal nucleation sequences that ex-situ measurements cannot reconstruct \cite{bocker2012in,mckenzie2020nucleation,mckenzie2021nucleation,aspillaga2023nucleation}. Spectroscopy and total-scattering methods probe local coordination in glasses and nanocrystalline products, where short-range order can persist after the long-range phase has changed \cite{mastelaro2018x,abdelhameed2023structural,shaaban2021spectroscopic,topper2023zinc}. Thermal analysis supplies onset temperatures but does not identify a phase without diffraction or chemical mapping. The strongest purity claim therefore requires at least two independent structural or compositional methods.

The phase-bottlenecks map (Figure~\ref{fig:phase-bottlenecks}) summarizes the
resulting decision logic: establish the actual composition, equilibrate or quench
at controlled points, identify every crystalline and glassy fraction, and only
then compare with phase-diagram constraints. This order prevents a visually
plausible diffraction match from being treated as a complete phase assignment.
